\chapter{Differentiation for arbitrary Euclidean Radon measures}\label{sec:radon-measure-differentiation}
For a singular Radon measure, a fivefold dilation of a ball may increase its mass without any fixed bound. The Euclidean differentiation theorem therefore needs a covering input stronger than the fivefold selection used for Lebesgue measure. We isolate that input and prove its analytic consequences. The source is \cite{Sim18}, Chapter 1, Lemma 4.5 and Corollary 4.10, printed pp.~21--24. In this section closed balls are used for the covering argument, and the final result is also stated for open balls.

\begin{theorem}[Besicovitch selection for Euclidean balls]\label{thm:besicovitch-ball-selection}
For each $n$ there is an integer $N_n$ with the following property. Let $\mathcal B$ be a collection of closed balls of positive radii in $\mathbb R^n$, with bounded centers and bounded radii, and let $A$ be their set of centers. There are at most $N_n$ finite or countable subcollections $\mathcal B_1,\ldots,\mathcal B_{N_n}$ such that each $\mathcal B_i$ consists of pairwise disjoint balls and
\[
 A\subset\bigcup_{i=1}^{N_n}\bigcup_{B\in\mathcal B_i}B.
\]
This is a cited geometric input: \cite{Sim18}, Chapter 1, Lemma 4.5, p.~21, where the proof is referred to \cite{Fed69}. It is not proved by the fivefold covering lemma, and its proof remains external here. The dimension-dependent bound is essential; a bound independent of the number of balls is being asserted.
\end{theorem}

\begin{proposition}[A finite-measure fine cover has an almost-everywhere disjoint selection]\label{prop:radon-symmetric-vitali}
Let $\lambda$ be a positive Radon measure on $\mathbb R^n$, let $A$ be a bounded Borel set with $\lambda(A)<\infty$, and let $\mathcal B$ be a family of closed balls centered in $A$ which covers $A$ finely: for each $x\in A$ there are balls in the family centered at $x$ with arbitrarily small positive radii. Then there is a finite or countable pairwise disjoint subfamily $\{B_j\}$ such that
\[
 \lambda\bigl(A\setminus\bigcup_jB_j\bigr)=0.
\]
This is the symmetric Vitali consequence proved in \cite{Sim18}, Remark 4.4(3), pp.~21--22, in the bounded-set form needed for local applications.
\end{proposition}
\begin{proof}
If $\lambda(A)=0$, the empty collection suffices. Otherwise restrict to radii at most one and apply Theorem \ref{thm:besicovitch-ball-selection}. Since the at most $N_n$ unions cover $A$, at least one union of disjoint balls covers a subset of $A$ of measure at least $\lambda(A)/N_n$. A finite subcollection of that union covers measure at least $\lambda(A)/(2N_n)$, by continuity from below. Keep those finitely many balls. Their union is closed.

For each point of $A$ outside this finite closed union, sufficiently small balls from its original fine cover miss the union. Apply the same argument to the remaining set using only such balls. The newly chosen finite family is disjoint both internally and from all previously chosen balls. At each stage the uncovered measure is multiplied by at most $1-1/(2N_n)$. If it becomes zero, stop; otherwise continue. The total selected family is countable, is pairwise disjoint, and leaves measure at most $(1-1/(2N_n))^k\lambda(A)$ after $k$ stages. Letting $k\to\infty$ proves the result. Closedness of each finite previous union is the reason the next fine cover remains available.
\end{proof}

\begin{proposition}[A maximal estimate without a doubling hypothesis]\label{prop:radon-maximal-weak}
For a positive Radon measure $\lambda$ on $\mathbb R^n$ and $g\in L^1(\lambda)$, put
\[
 M_\lambda g(x)=\sup_{0<r\leq1}
   \frac{\int_{\overline{B_r(x)}}|g|\,d\lambda}
        {\lambda(\overline{B_r(x)})},
\]
assigning value zero to a ratio with zero denominator. Then $M_\lambda g$ is Borel for a Borel representative of $g$, and for every $t>0$,
\[
 \lambda(\{M_\lambda g>t\})\leq\frac{N_n}{t}\|g\|_{L^1(\lambda)}.
\]
No doubling bound for $\lambda$ is assumed. This is a maximal-function formulation of the covering method behind \cite{Sim18}, Chapter 1, Section 4.
\end{proposition}
\begin{proof}
For a fixed radius, both ball integrals are Borel functions of the center by integration of the joint Borel ball indicator. On a compact ball both measures are finite. For fixed center, continuity from above of closed balls makes a positive-denominator ratio right continuous in the radius. Thus the displayed supremum can be taken over rational radii in $(0,1)$ together with radius one; zero-denominator terms cannot increase a nonnegative supremum. This proves measurability.

For each $x$ in the bounded set $\{M_\lambda g>t\}\cap B_R(0)$, take the family of all witness balls centered at $x$ for which the ratio is greater than $t$. Their radii are at most one. Theorem \ref{thm:besicovitch-ball-selection} selects at most $N_n$ disjoint subfamilies covering these centers. Every selected ball $B$ satisfies $t\lambda(B)<\int_B|g|\,d\lambda$. Sum this inequality within each disjoint family and then over the families. The measure of the covered centers is bounded above by the sum of the selected ball masses, while each family's integral sum is at most $\|g\|_1$. Therefore the desired bound holds on $B_R(0)$. Let $R\to\infty$ and use continuity from below. The argument does not enlarge a selected ball, so no comparison between its mass and a dilated ball's mass is needed.
\end{proof}

\begin{theorem}[Absolute oscillation vanishes at almost every Radon-measure point]\label{thm:radon-absolute-oscillation}
If $\lambda$ is a positive Radon measure on $\mathbb R^n$ and $g\in L^1(\lambda)$, then for $\lambda$-almost every $x$,
\[
 \lim_{r\downarrow0}\frac{1}{\lambda(\overline{B_r(x)})}
       \int_{\overline{B_r(x)}}|g(y)-g(x)|\,d\lambda(y)=0.
\]
This is an absolute-oscillation form of the differentiation conclusion in \cite{Sim18}, Corollary 4.10. Its covering input is the cited Besicovitch theorem above.
\end{theorem}
\begin{proof}
Almost every point lies in the support of $\lambda$ and has a finite value of $g$. At these points all the denominators are positive. Let $Dg(x)$ be the limsup on the left. It is measurable: the same joint-kernel argument and rational-radius reduction as in Proposition \ref{prop:radon-maximal-weak} apply to the locally integrable function $|g(y)-g(x)|$ for fixed $x$.

For $\phi\in C_c(\mathbb R^n)$, the triangle inequality and continuity at $x$ give
\[
 Dg(x)\leq M_\lambda(g-\phi)(x)+|g(x)-\phi(x)|.
\]
The average oscillation of $\phi$ tends to zero because it is bounded by its supremum oscillation on the shrinking ball, regardless of the measure of that ball. For $\varepsilon>0$, the maximal estimate and Chebyshev's inequality therefore give
\[
 \lambda(\{Dg>2\varepsilon\})
 \leq\frac{N_n+1}{\varepsilon}\|g-\phi\|_{L^1(\lambda)}.
\]
Theorem \ref{thm:radon-lp-density} supplies continuous compactly supported approximants with arbitrarily small right side. Hence this exceptional set has measure zero. Taking a countable sequence $\varepsilon\downarrow0$ proves $Dg=0$ almost everywhere. This proves the oscillation conclusion from the stated geometric covering input, rather than assuming a doubling property for a potentially singular measure.
\end{proof}

\begin{theorem}[Differentiation of functions relative to a Euclidean Radon measure]\label{thm:radon-function-differentiation}
Let $\lambda$ be a positive Radon measure on an open set $U\subset\mathbb R^n$, and let $f\in L^1_{\mathrm{loc}}(\lambda)$. For $\lambda$-almost every $x\in U$,
\[
 \lim_{r\downarrow0}\frac{1}{\lambda(B_r(x))}
       \int_{B_r(x)}|f(y)-f(x)|\,d\lambda(y)=0,
 \qquad
 \lim_{r\downarrow0}\frac{1}{\lambda(B_r(x))}
       \int_{B_r(x)}f\,d\lambda=f(x).
\]
The denominators are positive at almost every point for all sufficiently small radii. The average conclusion applies componentwise to a finite-dimensional vector-valued function. This is \cite{Sim18}, Chapter 1, Corollary 4.10, pp.~23--24, with the Euclidean covering input isolated in Theorem \ref{thm:besicovitch-ball-selection}.
\end{theorem}
\begin{proof}
First assume $U=\mathbb R^n$ and $f\in L^1(\lambda)$. Theorem \ref{thm:radon-absolute-oscillation} gives the closed-ball conclusion. At a point where it holds, fix a sufficiently small open-ball radius $r$ and take radii $s_k\uparrow r$. Their closed balls increase to $B_r(x)$. Continuity from below for the denominator and the nonnegative oscillation integral shows that every uniform bound for closed-ball averages at radii below a given scale also holds for open balls below that scale. Therefore the open-ball absolute oscillation tends to zero. The absolute value of the difference between the average of $f$ and $f(x)$ is bounded by that oscillation.

For the local statement, exhaust $U$ by relatively compact open sets $V_j$ and choose larger relatively compact open sets $W_j$ containing their closures. Restrict $\lambda$ to $W_j$ and extend this finite Borel measure by zero to $\mathbb R^n$; it is a Radon measure there. Restrict $f$ to $W_j$ and set it to zero elsewhere. Local integrability makes this a globally integrable function for the restricted measure. The whole-space result applies. At each point of $V_j$, all sufficiently small balls lie in $W_j$, so both integrals and denominators agree with the original ones. A countable union of exceptional null sets finishes the proof. Intersect the finitely many componentwise full-measure sets for a vector-valued function. The unproved geometric input remains Besicovitch selection, not the subsequent analytic differentiation argument.
\end{proof}
